% =====================================================================
% REPLACEMENT FOR SUPPLEMENTARY SECTION S15
%
% Every number marked \NUM{} is a PLACEHOLDER produced by the pipeline in
% scripts/. Run the three scripts, then substitute. Do not submit with
% \NUM{} markers present -- they render in bold red by design.
%
% \providecommand{\NUM}[1]{\textbf{\textcolor{red}{[#1]}}}   % add to preamble
% =====================================================================

\section*{S15. Perturbation as a model index: rapamycin and mTORC1}

Sections~S13.6 identifies biomedical agents as a prospective upstream layer.
A second is large-scale perturbation transcriptomics. Tahoe-100M is a
single-cell perturbation atlas of roughly 100 million transcriptomic profiles
measuring the response of 50 cancer cell lines to approximately 1{,}100
small-molecule perturbations (Zhang et al., 2025). This section shows how such
a resource enters Scientific Jigsaw, and --- more importantly --- the sharp
distinction between the two ways it can legitimately do so.

\subsection*{S15.1 Two roles for a perturbation, and the rule separating them}

A perturbation can enter a declared model in exactly two ways, and they are not
interchangeable.

\begin{enumerate}
\item \textbf{As a model index} (Section~S11.11). When independent structural or
biochemical evidence establishes that a drug creates or destroys a molecular
relation, the treated and untreated conditions are \emph{different declared
models} with \emph{different admissible spaces}:
\[
\Omega_{\mathrm{control}} \;\neq\; \Omega_{\mathrm{drug}}.
\]
\item \textbf{As an external weight} (Section~S11.12). When the perturbation
changes only relative component abundance, the admissible set is unchanged and
only the distribution over it moves:
\[
\Omega_{\mathrm{control}} = \Omega_{\mathrm{drug}}, \qquad
p(\omega \mid \mathrm{control}) \neq p(\omega \mid \mathrm{drug}).
\]
\end{enumerate}

\noindent In neither case does the transcriptional response itself supply a
contact, a prerequisite or an exclusion. Co-regulation is not co-assembly, and a
differential-expression statistic is not a temporal claim. What the atlas
supplies is the \emph{condition label} that selects a declared model variant,
and the \emph{component availability} that fixes the board's component set.
Everything mechanistic is declared separately and cited separately.

\subsection*{S15.2 The rapamycin model index}

Rapamycin is the clearest available case of role~1 because its mechanism is a
molecular assembly event, established independently of any transcriptomic
measurement. Rapamycin binds FKBP12, and the FKBP12--rapamycin complex binds the
FRB domain of mTOR (PDB 1FAP; and within the complex, PDB 5FLC). Binding
additionally \emph{disrupts} the mTOR--RAPTOR association and the mTORC1 dimer
--- a result from the biochemical and cryo-EM literature, not from any single
deposited structure, and cited as such.

Two declared model variants follow over the mTORC1 board
$\{\mathrm{mTOR}, \mathrm{RAPTOR}, \mathrm{mLST8}, \mathrm{PRAS40},
\mathrm{DEPTOR}\}$:

\begin{itemize}
\item $M_0$ (untreated): contacts extracted from apo mTORC1 coordinates
(PDB 6BCX; cross-checked against 6BCU) under the criterion used throughout ---
heavy-atom separation within 5~\AA, minimum three unique interface residues.
\item $M_{\mathrm{Rapa}}$ (rapamycin): $M_0$ plus one declared
\emph{contact}, FKBP12--mTOR(FRB), conditional on rapamycin; minus one declared
\emph{exclusion}, mTOR--RAPTOR.
\end{itemize}

\noindent This is the first case in the manuscript in which a single declared
change exercises two relation types at once, and the first in which an exclusion
is the mechanistically decisive relation.

\begin{table}[h]
\centering\small
\begin{tabular}{lrrl}
\hline
model & components & merger trees & fragments \\
\hline
$M_0$ (untreated)          & 5 & \NUM{n0} & \NUM{frag0} \\
$M_{\mathrm{Rapa}}$        & 6 & \NUM{nR} & \NUM{fragR} \\
\hline
\end{tabular}
\caption{Exact consequence of the declared rapamycin relations. The perturbation
does not reweight the untreated space: it changes which histories exist.}
\label{tab:rapaspaces}
\end{table}

The mechanistically informative output is not the change in count but the
change in \emph{connectivity}. The declared exclusion severs mTOR from RAPTOR,
so the contact graph of the treated model is disconnected and the holo-complex
is not merely less probable but \textbf{impossible}; the space resolves into
independently formable fragments. Because the merger-tree enumerator requires a
connected contact graph (Section~S11.5), the framework reports this as an exact
structural verdict and names the excluded relation responsible. This recovers
the known biochemistry --- rapamycin releases RAPTOR from the mTOR core --- as a
consequence of the declared model rather than as an input to it.

\subsection*{S15.3 Cellular context as a second axis}

Tahoe's design is a grid of compounds against cell lines, and Scientific Jigsaw
consumes both axes. The compound indexes the model variant, as above. The cell
line fixes the component set: a subunit not expressed in a given line is not
available for assembly and is removed from the board, which changes the
admissible histories rather than their weights.

The evidential asymmetry here is important and works in the framework's favour.
Absence of transcript is strong evidence --- no transcript implies no protein,
hence no component. Presence is weak --- transcript may not yield
assembly-competent protein. Verdicts of the form \emph{impossible in this
context} are therefore robust, and verdicts of the form \emph{formable} are
permissive only. Exact enumeration delivers impossibility most confidently, so
the negative calls are reported and the positive ones are not treated as
predictions.

\begin{table}[h]
\centering\small
\begin{tabular}{lrrl}
\hline
context & $M_0$ trees & $M_{\mathrm{Rapa}}$ trees & mTOR:RAPTOR:mLST8 \\
\hline
reference (all subunits expressed) & \NUM{c0} & \NUM{cR} & \NUM{v0} \\
\NUM{line-specific rows from tahoe\_availability.py} & & & \\
\hline
\end{tabular}
\caption{Perturbation and cellular context act on the admissible space
independently and exactly. Rows are generated by the pipeline in the
reproducibility release; the expression threshold is swept as in Section~S5.2.}
\label{tab:rapacontexts}
\end{table}

\subsection*{S15.4 The contrasting null: coordinated regulation}

The proteasome provides the contrasting case, and it is instructive precisely
because the answer is that nothing happens. Proteasome inhibition triggers the
NRF1 (NFE2L1)-mediated ``bounce-back'', in which proteasome subunit genes are
\emph{coordinately} upregulated (Radhakrishnan et al., 2010). Under the
encounter-order prior of Section~S5.9, a uniform fold-change is exactly
invariant: replacing every weight $w_x$ by $cw_x$ leaves
$w_{x_k} / \sum_{j \in R_k} w_j$ unchanged at every step, so every history
probability, the entropy and every subcomplex support are identical to control.

\begin{table}[h]
\centering\small
\begin{tabular}{lrrrr}
\hline
condition & $N_{\mathrm{eff}}$ & $13\mathrm{S}{+}\beta_1$
 & $13\mathrm{S}{+}\beta_5$ & $13\mathrm{S}{+}\beta_1{+}\beta_5$ \\
\hline
control $(1,1,1)$                & 6.00 & 0.333 & 0.333 & 0.333 \\
coordinated $\times 3$ $(3,3,3)$ & 6.00 & 0.333 & 0.333 & 0.333 \\
differential $(1,3,1)$           & 4.91 & 0.200 & 0.600 & 0.450 \\
\hline
\end{tabular}
\caption{Synthetic illustration on the six admissible 20S orders of
Section~S14.2. Coordinated upregulation --- the response actually elicited by
proteasome inhibition --- is the exact null: only a change in \emph{relative}
subunit availability moves the distribution. The admissible set is six orders in
every row. These values are invented for exposition and are not measurements.}
\label{tab:weightnull}
\end{table}

Taken together the two cases delimit what perturbational data can do. Rapamycin
changes the model; coordinated regulation changes nothing; only differential
availability reweights. Reporting a reweighted support as though a route had
been excluded --- or a coordinated fold-change as though it had any effect at
all --- would be precisely the conflation the framework exists to prevent.

\subsection*{S15.5 Scope and what would strengthen this}

The rapamycin relations are declared, not inferred, and the analysis inherits
their status: the FKBP12--FRB contact is structurally grounded, whereas the
mTOR--RAPTOR exclusion is a declared idealisation of a disruption reported as
weakening rather than abolition, and is queried with that qualification. The
board treats the mTORC1 monomer and does not represent the dimer, so
dimer-specific consequences are outside its scope. Transcript abundance is a
weak proxy for assembly-competent protein, dose and exposure time condition
every availability call, and the analysis is unweighted throughout.
Reciprocal validation --- testing whether rapamycin-insensitive lines are those
in which the declared exclusion leaves the observed intermediates formable ---
remains the natural next test and is not attempted here.
